\documentclass[12pt,a4paper]{article}
\usepackage[utf8]{inputenc}
\usepackage[vietnamese]{babel}
\usepackage[left=2cm,right=2cm,top=2cm,bottom=2cm]{geometry}
\usepackage{mathptmx}
\usepackage{enumitem}

\pagestyle{plain}
\linespread{1.15}

\begin{document}

% --- HEADER QUỐC HIỆU & CƠ QUAN BAN HÀNH ---
\noindent
\begin{minipage}[t]{0.45\textwidth}
    \begin{center}
        \small
        CÔNG AN THÀNH PHỐ HÀ NỘI\\
        \textbf{\underline{PHÒNG CẢNH SÁT HÌNH SỰ}}\\[0.1cm]
        Số: 11c/BB-HC\\
        \textit{V/v hỏi cung lấy lời khai tự thú}
    \end{center}
\end{minipage}
\hfill
\begin{minipage}[t]{0.52\textwidth}
    \begin{center}
        \small
        \textbf{CỘNG HÒA XÃ HỘI CHỦ NGHĨA VIỆT NAM}\\
        \textbf{\underline{Độc lập - Tự do - Hạnh phúc}}\\[0.1cm]
        \textit{Hà Nội, ngày 25 tháng 07 năm 2026}
    \end{center}
\end{minipage}

\vspace{0.6cm}

% --- TIÊU ĐỀ VĂN BẢN ---
\begin{center}
    \textbf{\large BIÊN BẢN HỎI CUNG BỊ CAN / TỰ THÚ TOÀN BỘ}\\
    \textit{(Đối tượng Nguyễn Thanh Tùng thừa nhận hành vi xô xát đêm 24/07)}
\end{center}

\vspace{0.3cm}

\noindent Vào hồi 19 giờ 00 phút ngày 25 tháng 07 năm 2026, tại Phòng Hỏi cung số 1, Trụ sở Phòng CSHS.\\
Chúng tôi gồm:
\begin{enumerate}[label=\arabic*., leftmargin=1.5cm]
    \item Điều tra viên: Đại úy Lê Minh — Đội Trọng án.
    \item Cán bộ ghi biên bản: Trung úy Nguyễn Văn Hoàng — Cán bộ điều tra.
\end{enumerate}

\noindent Tiến hành hỏi cung đối tượng: **NGUYỄN THANH TÙNG (Sinh năm 1990)**.\\
Trước các chứng cứ vật chất đanh thép: Chiếc ghim cài áo rơi tại hiện trường (`p8`), Khung tranh thu giữ tại phòng trọ (`f3-2`), và Mảnh bài báo năm 1998 (`p5`), Nguyễn Thanh Tùng đã khai nhận toàn bộ sự thật:

\section*{NỘI DUNG TỰ THÚ (HỎI VÀ ĐÁP)}

\noindent \textbf{Hỏi (ĐTV Lê Minh):} Anh hãy tường trình lại nguồn cơn mâu thuẫn và diễn biến cuộc chạm trán giữa anh và Khang tối 24/07?\\
\textbf{Đáp (Nguyễn Thanh Tùng):} *(Cúi gục đầu, run rẩy bật khóc)* Thưa cán bộ, hai hôm trước ngày xảy ra vụ án, mấy đứa bạn cũ cùng xóm có rủ tôi đi nhậu chia tay trước khi tôi về công trường Hải Phòng. Trong lúc say ngà ngà, cả nhóm ôn lại chuyện nghịch ngợm hồi nhỏ. Thằng Khang phấn khích cười phá lên bảo: \textit{"Hồi đó tao trốn đỉnh nhất là vụ tao gài chốt nhốt thằng nhóc vào tủ, tụi mày tìm cả ngày không ai nghĩ ra luôn!"}.

Lúc đó tôi chết lặng giữa bàn nhậu... Tôi bàng hoàng nhận ra suốt 12 năm qua, em trai Gia Huy của tôi không phải tự chui vào tủ bị tai nạn, mà chính là do thằng Khang đã cố tình gài chốt nhốt em tôi đến chết ngạt!

\noindent \textbf{Hỏi (ĐTV Lê Minh):} Tối ngày 24/07, anh đã làm gì?\\
\textbf{Đáp (Nguyễn Thanh Tùng):} Tôi không chịu nổi sự thật đó. Tầm 20:00 tối 24/07, tôi cầm theo bài báo cũ năm 1998 sang nhà Khang chất vấn. Tôi hỏi nó tại sao năm xưa làm chết em tôi mà lại nói dối trước báo chí. Thằng Khang nghe xong chẳng những không hối lỗi mà còn cười khẩy bảo: \textit{"Thôi mà, chuyện trẻ con nghịch ngợm xưa như trái đất rồi nhắc lại làm gì!"}, rồi nó giằng lấy tờ báo trên tay tôi xé vụn vứt xuống sàn thách thức!

Quá uất ức và điên tiết, tôi lao vào giằng co với Khang. Hai bên xô đẩy làm gạt đổ vỡ toang bộ bình trà trên bàn xuống sàn gạch. Tôi dùng hết sức xô mạnh một cái làm Khang ngã ngửa đập gáy vào mép tủ âm tường rồi ngất xỉu nằm bất động.

\noindent \textbf{Hỏi (ĐTV Lê Minh):} Sau khi Khang ngã ngất, anh đã làm gì tiếp theo?\\
\textbf{Đáp (Nguyễn Thanh Tùng):} Tôi thấy Khang nằm im không nhúc nhích thì hoảng sợ tột cùng. Tôi quỳ xuống lay người và áp tai vào ngực kiểm tra thì thấy **Khang vẫn còn thở và mạch còn đập**. Vì quá sợ bị liên lụy nên lúc 20:15 tôi chạy xộc ra ngõ, phóng xe máy ra bến đón xe khách về Hải Phòng. 

Lúc ra bến xe, tôi hoang mang không biết Khang đã tỉnh lại chưa. Đến 20:55, tôi bấm số gọi vào máy Khang để thăm dò: nếu Khang nghe máy thì nghĩa là Khang không sao, còn nếu công an có hỏi thì tôi sẽ bảo chỉ gọi điện hỏi vay vốn làm ăn để chối tội. Tôi thề có linh hồn em trai tôi, tôi chỉ xô ngã Khang chứ lúc tôi rời đi Khang vẫn còn sống, tôi hoàn toàn không hề đâm chết Khang!

\vspace{0.6cm}
\noindent Biên bản kết thúc hồi 20 giờ 30 phút cùng ngày. Nguyễn Thanh Tùng đã đọc lại và ký xác nhận.

\vspace{0.8cm}

% --- CHỮ KÝ NGHỊ ĐỊNH 30 ---
\noindent
\begin{minipage}[t]{0.45\textwidth}
    \small
    \textbf{\textit{Nơi nhận:}}\\
    - Viện KSMND TP. Hà Nội;\\
    - Hồ sơ chuyên án TEST-99;\\
    - Lưu: VT, CSHS.
\end{minipage}
\hfill
\begin{minipage}[t]{0.5\textwidth}
    \begin{center}
        \small
        \textbf{NGƯỜI TỰ THÚ}\\
        \textit{(Ký, ghi rõ họ tên)}\\[1.5cm]
        \textbf{Nguyễn Thanh Tùng}
    \end{center}
\end{minipage}

\end{document}
